% =============================================================================
% Section 9: Conclusions
% =============================================================================

\section{Conclusions}
\label{sec:conclusions}

This study developed and verified a closed, physics-based dimensionless framework for predicting bypass flow partitioning, convective heat transfer, and pumping penalties in single-phase liquid immersion cooling of server electronics. The key conclusions are:

\begin{enumerate}
\item \textbf{Physics-Based Flow Partitioning}: The internal bypass split is governed by the parallel hydraulic resistance ratio between the active fin matrix and overhead clearance passages. The derived dimensionless closure predicts bypass mass fraction ($\Phi_{\mathrm{bypass}}$) across $\mathrm{OR} \in [0.0, 1.0]$ and $\mathrm{Re}_{\mathrm{ch}} \in [25, 1000]$ with a Mean Absolute Error of $5.72$~percentage points and an unsealed relative percentage error of $\mathbf{13.97\%}$ ($\mathrm{OR} \ge 0.10$) on the calibration dataset.
\item \textbf{Thermally Developing Convective Heat Transfer}: Coupling the partitioned active Reynolds number $\mathrm{Re}_{\mathrm{active}}$ with an asymptotic composite Graetz model captures thermal entrance development across high Prandtl number dielectric coolants ($\mathrm{Pr} = 9.2\text{--}210$), yielding thermal resistance predictions within $\mathbf{0.75\%}$ MAPE ($R^2 = 0.9722$).
\item \textbf{Out-of-Sample Generalizability}: When evaluated on withheld dielectric coolants (EFL-1 and PAO-4) and distinct heat-sink geometries (Micro-Pin-Fin and Oblique-Fin) without refitting, the dimensionless framework demonstrated robust predictive generalizability with out-of-sample thermal resistance errors bounded below $\mathbf{0.84\%}$ and bypass MAE within $6.46$~percentage points. For high-Prandtl oil (PAO-4, $\mathrm{Pr}=72.0$), the uncalibrated Nusselt prediction exhibits a $31.79\%$ error due to boundary-layer viscosity gradients, highlighting the physical role of fluid property variations.
\item \textbf{Thermal-Hydraulic Optimization}: Evaluating the global Figure of Merit reveals that structural open ratios $\mathrm{OR} \le 0.15$ provide optimal cooling efficiency, realizing over $90\%$ of the maximum achievable thermal enhancement while avoiding excessive hydraulic pumping power consumption.
\item \textbf{Verification and Validation Rigor}: The OpenFOAM numerical framework was verified against a multi-tier reproducibility ladder across six literature benchmarks (Kewalramani et al., Sastre et al., Dong et al., Chun et al., Khoshvaght-Aliabadi et al., and Jeng), alongside a four-grid E\c{c}a--Hoekstra / Roache GCI study achieving fine-grid numerical uncertainties below $0.70\%$ for temperature and $0.64\%$ for bypass fraction.
\end{enumerate}
